\section{Divisorial corrections and surface links}
\label{sl:section}

We first remove the contributions vertical over the common base of a link.
When that base has dimension two, the remaining surface calculation
produces a correction depending on its function field;
Section~\ref{sf:section} will make the correction intrinsic to each
Mori fibre space for uniformly large test degree.

\subsection{Vertical divisors and change of base}

Fix a degree $d$ and a subcollection $\mathcal S$ of the Picard-number-one
K3 surfaces of degree $d$ with trivial automorphism group.  Write $u(Y)$ for
the indicator that the maximal rationally connected quotient of a smooth
proper model of $\C(Y)$ is birational to a member of $\mathcal S$.
As in Definition~\ref{pre:test}, this depends only on the function
field, and $u$ is extended additively to finite lists of fields.  All
fourfolds in this section are the terminal Mori fibre spaces and
intermediate models in the ordinary Sarkisov paths supplied by
Theorem~\ref{bd:diagrams}, birational to the fixed rationally connected
fourfold.  Every reference to an ordinary link below is to a link in
one of these chosen paths.  In the eventual application that
fourfold is assumed rational.  The common base of a link is denoted by
$Q$.

For a proper contraction $f:Y\to Z$ with the generic-fibre property in
Lemma~\ref{pre:mrc-fibration}, let $\operatorname{Div}_{Z}(Y)$ denote the
prime divisors on $Y$ whose images are proper subsets of $Z$.
We define
\begin{equation}
\label{sl:v-definition}
 v(Y/Z)=
 \sum_{E\in\operatorname{Div}_{Z}(Y)}u(E)
 -\sum_{D\in\operatorname{Div}(Z)}u(D).
\end{equation}
This notation means a finite cancelled sum, as follows.  There is a
dense open $U\subset Z$ such that, for every prime divisor $D$ of $Z$
meeting $U$, precisely one prime divisor $E_D$ of $Y$ dominates $D$,
and $E_D\dashrightarrow D$ has the generic-fibre property of
Lemma~\ref{pre:mrc-fibration}.  To obtain $U$, spread a resolution of the
generic fibre, and shrink so that the resulting smooth proper family
has geometrically integral rationally connected fibres and the
birational map to the original family remains fibrewise birational
over dense opens.  Generic flatness and openness of geometric
integrality give uniqueness of $E_D$ after a further shrinking.
Choose the same open so that $f$ is flat over $U$.  By the dimension
formula, a prime divisor meeting $f^{-1}(U)$ and vertical over $U$
must dominate a divisor of $U$.  It is therefore one of the divisors
$E_D$ just described.
Cancel $u(E_D)-u(D)=0$ for every such $D$.
Every remaining prime divisor of $Z$ is a component of $Z\setminus U$,
and every remaining prime divisor of $Y$ is a component of
$f^{-1}(Z\setminus U)$.  Both are finite lists.  Replacing $U$ by a
smaller good open only adds pairs of equal terms, so the resulting
integer is independent of $U$.

\begin{lemma}[Uniruled divisors in contraction fibres]
\label{sl:uniruled-vertical}
Let $f:Y\to Z$ be a projective contraction of normal varieties over an
algebraically closed field of characteristic zero.  Suppose that $Y$
is of Fano type over $Z$.  Let $E$ be a prime divisor which is a
component of $f^{-1}(f(E))$ and satisfies
$\dim E>\dim f(E)$.  Then $E$ is uniruled.
Consequently:
\begin{enumerate}[label=(\roman*)]
\item if $\dim Y>\dim Z$, every divisor vertical over $Z$ is uniruled;
\item if $f$ is birational, every $f$-exceptional prime divisor is
uniruled.
\end{enumerate}
The same conclusions apply to geometric generic fibres of such a
contraction and to birational transforms of the divisors.
\end{lemma}

\begin{proof}
Choose a rational boundary $\Theta$ making $(Y,\Theta)$ klt and
$-(K_Y+\Theta)$ relatively ample. Then
$-K_Y=-(K_Y+\Theta)+\Theta$ is relatively big, and a sufficiently divisible
multiple of $-(K_Y+\Theta)$ is generated over $Z$.
Theorem~1.2 and Corollary~1.4 of \cite{HM07Vertical} give rational chain
connectedness of every fibre; the non-klt locus is empty.
The statement after passage to a characteristic-zero geometric generic
field follows by descent of the finite data and base extension.

Put $T=f(E)$ and choose a general point $e$ of $E$ outside the other
components of $f^{-1}(T)$.  The fibre of $E\to T$ through $e$ has
positive dimension.  A chain of rational curves in the fibre of $f$
joining $e$ to another point starts with a nonconstant rational curve
through $e$.  That curve is contained in $E$: its image is in
$f^{-1}(T)$, and an irreducible curve through a point lying on only the
component $E$ must lie in that component.  Thus rational curves cover
$E$, proving uniruledness.  Equivalently, one can select a dominating
component of the parameter space of these rational curves.  The
argument can be made after an uncountable algebraically closed
extension and descended, so it does not require an uncountability
assumption on the original field.

For (i), $f(E)$ is proper and $f^{-1}(f(E))$ is a proper closed subset
of the integral variety $Y$.  Its component containing $E$ must equal
$E$, and
$\dim E-\dim f(E)\ge\dim Y-\dim Z>0$.  The same component argument
works in (ii), where exceptionalness gives
$\dim E-\dim f(E)\ge1$.  Finally uniruledness is birationally
invariant.
\end{proof}

The contractions $B\to Q$ which occur for the bases of link ends are
of Fano type over $Q$, by the base-contraction part of
Theorem~\ref{bd:diagrams}.  Their geometric generic fibres are
geometrically integral and rationally connected; this also holds for
$V\to Q$ by composition.  Thus all the following cancelled sums are
defined.

For a tower $V\to B\to Q$, let $h_Q(V/B)$ be the portion of
\eqref{sl:v-definition} consisting of divisors dominating $Q$:
\begin{equation}
\label{sl:h-definition}
 h_Q(V/B)=
 \sum_{\substack{E\in\operatorname{Div}_{B}(V)\\E\to Q\ {
 dominant}}}u(E)
 -\sum_{\substack{D\in\operatorname{Div}(B)\\D\to Q\ {
 dominant}}}u(D).
\end{equation}
Use the same cancellation as for $v(V/B)$.  A cancelled pair $E_D,D$
has the same dominance behaviour over $Q$, so this is well defined.
Partitioning both finite cancelled lists according to dominance over
$Q$ gives
\begin{equation}
\label{sl:tower}
 v(V/Q)=v(V/B)-h_Q(V/B)+v(B/Q).
\end{equation}
Indeed a prime divisor of $V$ vertical over $Q$ is necessarily
vertical over $B$.  In the right-hand side, the divisors of $B$
vertical over $Q$ cancel between the first two terms and $v(B/Q)$,
leaving exactly the defining sum for $v(V/Q)$.  Good opens may be
shrunk simultaneously to perform these cancellations on finite lists.

\begin{lemma}
\label{sl:base-v-zero}
For every base contraction $B\to Q$ of an ordinary fourfold link,
$v(B/Q)=0$.
\end{lemma}

\begin{proof}
The assertion is immediate for an isomorphism.  For a birational
contraction, nonexceptional prime divisors cancel with their
birational transforms on $Q$.  The remaining divisors are uniruled by
Lemma~\ref{sl:uniruled-vertical}.  Since $\dim B\le3$, these divisors
have dimension at most two, and their MRC quotients have dimension at
most one; they have test value zero.

If the contraction has positive relative dimension, then
$\dim Q\le2$.  Every prime divisor of $Q$ has test value zero for
dimension reasons.  Every divisor of $B$ vertical over $Q$ is
uniruled by Lemma~\ref{sl:uniruled-vertical}, and again has dimension
at most two, so its test value is zero.  This includes $Q$ a point,
when there are no vertical divisors and no base divisors to sum.
\end{proof}

For a link $f:V/B\dashrightarrow V'/B'$ over $Q$, split its divisorial
cocycle into the contributions from valuations vertical and horizontal
over $Q$, as in \cite[Definition~2.3]{LinShinder26}:
\[
 c_u(f)=c_{u,\mathrm{vert}/Q}(f)+c_{u,\mathrm{hor}/Q}(f).
\]
The signs are those of Lemma~\ref{pre:cocycle}: a divisor present
on the primed model and absent on the unprimed model contributes
positively.  A common prime divisor has the same residue field and the
same dominance behaviour over $Q$ on both models.

\begin{proposition}[Vertical correction]
\label{sl:vertical}
For every ordinary fourfold link over $Q$ one has
\begin{align}
\label{sl:vertical-cocycle}
 c_{u,\mathrm{vert}/Q}(f)&=v(V'/Q)-v(V/Q),\\
\label{sl:horizontal-residual}
 c_u(f)-\bigl(v(V'/B')-v(V/B)\bigr)
 &=c_{u,\mathrm{hor}/Q}(f)-h_Q(V'/B')+h_Q(V/B).
\end{align}
These identities hold for every test degree and require no
boundedness threshold.
\end{proposition}

\begin{proof}
In the difference $v(V'/Q)-v(V/Q)$ the base-divisor sums cancel.
Common vertical prime valuations cancel as well, leaving precisely
the signed exceptional divisors vertical over $Q$.  This proves
\eqref{sl:vertical-cocycle}.  Substitute \eqref{sl:tower} and use
Lemma~\ref{sl:base-v-zero} on both ends to obtain
\eqref{sl:horizontal-residual}.
\end{proof}

\subsection{The common-base dimensions other than two}

Denote the right-hand side of \eqref{sl:horizontal-residual} by
$H_Q(f)$.  We next specify exactly where a degree threshold is needed.

\begin{proposition}
\label{sl:low-base}
There is an integer $d_{\mathrm{vert}}$, depending only on the bounded
marked families in Theorem~\ref{bd:diagrams} and the exclusions in
Theorem~\ref{bd:exclusions}, with the following property.  For every
$d>d_{\mathrm{vert}}$, every allowed subcollection $\mathcal S$, and
every ordinary fourfold link $f$ whose common base satisfies
$\dim Q\le1$, one has
\[
 c_{u,\mathrm{hor}/Q}(f)=h_Q(V/B)=h_Q(V'/B')=0.
\]
In particular $H_Q(f)=0$.  If $\dim Q=3$, these equalities hold for
every $d$.  The integer $d_{\mathrm{vert}}$ is independent of the link,
the birational map, the number of links in a path, and the choice of
$\mathcal S$.
\end{proposition}

\begin{proof}
Suppose first that $\dim Q=3$.  Each Mori base has dimension at most
three and dominates $Q$, so its morphism to $Q$ is birational.
No divisor of either Mori base, and no divisor of $V$ vertical over
that base, dominates $Q$.  Therefore both $h_Q$ terms vanish.
Normal proper generic curves over $\C(Q)$ are regular and hence
smooth in characteristic zero.  A birational map between them is an
isomorphism.  Their closed points are exactly the prime divisors of
the fourfolds dominating $Q$, so there are no horizontal divisorial
contributions either.

For $\dim Q=0$, the entire marked contraction diagram is bounded by
Theorem~\ref{bd:diagrams}.  The exceptional prime divisors contributing
to $c_u(f)$ therefore have bounded birational models of dimension
three.  Choose a marked good open for each contraction $V\to B$.
After cancellation, the divisors contributing to $h_Q(V/B)=v(V/B)$
are components of the marked bad locus on $B$ or its preimage on
$V$.  These components have bounded birational models by the
\emph{marked-locus} assertion of Theorem~\ref{bd:diagrams}.
The bounded-variety exclusion in Theorem~\ref{bd:exclusions}\textup{(i)} gives a
single threshold beyond which every term has test value zero.

Now suppose $\dim Q=1$, and put $F=\C(Q)$.  We use the bounded
diagrams after extension to $\overline F$.  If $E$ is a horizontal
exceptional prime divisor of the link, the geometric generic fibre
components of $E\to Q$ are surfaces.  On the side where the divisor
is extracted, each is exceptional for a divisorial contraction of
the generic threefold.  They are uniruled by
Lemma~\ref{sl:uniruled-vertical}; any intervening small modifications
preserve their birational types.  Their smooth birational models are
bounded by Theorem~\ref{bd:diagrams}.  The threefold-over-a-curve
exclusion in Theorem~\ref{bd:exclusions}\textup{(iv)} therefore makes
$u(E)=0$ uniformly for large $d$.  Here $u$ is evaluated on the original
complex threefold field $\C(E)$; only its geometric generic surface
components, rather than $E$ itself, are being placed in bounded families.

It remains to treat $h_Q(V/B)$, the other end being identical.  Work
with a good open for the generic contraction
$V_F\to B_F$ and spread it over an open of $Q$.  Shrinking $Q$ changes
no horizontal prime field.  The cancellations defining $h_Q$ remove
all divisors meeting that good open.  Consequently the geometric
generic components of the remaining prime divisors lie in the
bounded marked bad locus or its bounded preimage.

A remaining positive term is a divisor $P\subset V$ vertical over
$B$ and dominant over $Q$.  Its geometric generic fibre components
are vertical divisors of $V_{\overline F}\to B_{\overline F}$.
They are uniruled surfaces by Lemma~\ref{sl:uniruled-vertical}, with
bounded birational models.  Theorem~\ref{bd:exclusions}\textup{(iv)}
therefore gives $u(P)=0$ for large $d$.  A remaining negative term is a divisor
$D\subset B$ dominant over $Q$.  If $\dim D\le1$, its test value is
zero automatically.  If $\dim D=2$, its geometric generic fibre
components over $Q$ are bounded curves.  Were $u(D)=1$, the surface
$D$ would be birational to a member of $\mathcal S$; the induced
pencil, after Stein factorization, would have general fibre genus at
least $1+\sqrt d/2$.  This contradicts the uniform genus bound for
those components when $d$ is large.

Only finitely many bounded families and types of marked components
have been used.  Take the maximum of their exclusion thresholds.
All exclusions concern the full collection of degree-$d$
Picard-number-one K3 surfaces and therefore hold for every
subcollection $\mathcal S$.  This proves the stated uniformity.
\end{proof}

\subsection{Reduction over a surface common base}

\begin{lemma}
\label{sl:generic-surface-reduction}
Suppose $\dim Q=2$.  Then $k=\C(Q)$ is a rational function field in
two variables over $\C$.  The generic fourfold-link models are smooth
projective geometrically integral surfaces over $k$, and all small
modifications in the link become isomorphisms.  An end is either a
del Pezzo surface of Picard number one over $k$, or a conic bundle
over a smooth projective geometrically integral genus-zero curve
over $k$, of relative Picard number one.  These statements do not
assert that the surface or the genus-zero curve is split over $k$.
In the conic case the total Picard number of the generic surface over
$k$ is two.
\end{lemma}

\begin{proof}
A smooth projective model of $Q$ is a dominant rational image of the
rationally connected fourfold, hence is rationally connected.  A
smooth projective rationally connected complex surface is rational:
it is uniruled, its minimal model is ruled over a curve, and rational
connectedness forces that curve to have genus zero; see
\cite[Theorem~30]{AraujoKollar02} and \cite[Theorem~V.1]{Beauville78}.  Consequently
$k\simeq\C(x,y)$.

A terminal fourfold is nonsingular in codimension two
\cite[Exercise~4.9.1.1]{KollarFlips057}, so its
singular locus has dimension at most one and cannot dominate $Q$.
The generic surface is regular, as a localization of its regular
total space near the generic fibre, and is smooth because $k$ is
perfect.  The contractions have connected fibres and rationally
connected geometric generic fibres as above, giving geometric
integrality.  Projectivity is inherited by base change.  A small
birational map between smooth proper surfaces is an isomorphism:
the exceptional locus of a non-isomorphic birational morphism of
smooth surfaces contains a curve, and resolving a birational map
gives the same conclusion if it is non-isomorphic in either
direction.  Such a curve would spread to a divisor exceptional for
the original small map.

If $\dim B=2$, the contraction $B\to Q$ is birational.  The generic
end $V_k$ has ample anticanonical divisor and Picard number one over
$k$, hence is a del Pezzo surface.  If $\dim B=3$, the generic base
$B_k$ is a normal proper curve with geometrically integral
rationally connected generic model, so it is a smooth genus-zero
curve.  The generic morphism $V_k\to B_k$ has relatively ample
anticanonical divisor and relative Picard number one, hence is a
smooth-surface Mori conic bundle.

Here Picard numbers mean dimensions of numerical divisor spaces
\emph{over $k$}, not over its algebraic closure.  To justify their
restriction from the fourfold, a divisor on $V_k$ closes up to a
Weil divisor on $V$, which is $\Q$-Cartier.  This gives surjectivity
onto the generic numerical divisor space.  A class numerically zero
over the indicated original base remains numerically zero on the
generic fibre: a generic curve and its intersection degrees spread
over an open of $Q$.  Thus the relative rank is at most the original
rank one.  The anticanonical class is nonzero, so it is exactly one.
The same argument bounds the numerical rank of the generic middle
surface by two.  No passage to an algebraic closure is made in this
rank calculation.
In the conic case the base curve has numerical Picard number one,
its pullback to the surface is nonzero, and the relative numerical
divisor space has dimension one.  Thus the total numerical divisor
space has dimension two.
\end{proof}

\begin{lemma}[The generic dichotomy]\label{sl:generic-dichotomy}
Suppose $\dim Q=2$ and put $k=\C(Q)$.
After identifying the middle generic surfaces by the restricted small
maps, let $W$ be their common smooth model.  Exactly one of the following
reductions applies:
\begin{enumerate}[label=\textup{(\roman*)}]
\item The restricted link induces an isomorphism of the structured
generic ends.  In the conic case their curve fields are the same
embedded subfield of $k(W)$.  Its horizontal cocycle and $H_Q$ vanish.
\item The two ends are obtained from $W$ by distinct elementary
$K_W$-negative contractions, each either divisorial to a rank-one del
Pezzo surface or a conic contraction.  Here $\rho(W/k)=2$, the two
contractions give the two extremal rays, and $-K_W$ is ample.
\end{enumerate}
In particular two distinct conic structures on the same surface belong
to \textup{(ii)}, not \textup{(i)}.
\end{lemma}

\begin{proof}
By Lemma~\ref{sl:generic-surface-reduction}, $\rho(W/k)\leq2$.
If a divisorial contraction $W\to S$ survives on the generic fibre,
then $\rho(W/k)=\rho(S/k)+1$.  The end $S$ therefore has rank one and
is a del Pezzo surface.  Extraction from a conic end, which has rank
two, would force $\rho(W/k)\geq3$.  Thus any surviving extraction forces
$\rho(W/k)=2$.  If the opposite end has no surviving extraction, its
surface is $W$ itself and cannot be a rank-one del Pezzo.  It must
instead carry its conic contraction.  If both ends have surviving
extractions, both are divisorial contractions from $W$ to rank-one
del Pezzo surfaces.

If neither side has a surviving extraction, both generic end surfaces
are $W$.  Rank one then gives two structures over a point, and the
restricted map is a structured isomorphism.  Rank two gives two conic
contractions.  These may agree or may be distinct; the latter case is
retained as a two-conic move.

In all the rank-two cases, compare the two end contractions on $W$.
If they contract the same extremal ray, uniqueness of the contraction
identifies their normal targets.  Indeed, each is constant on the
connected fibres of the other, hence factors through it, and the two
factorizations are inverse.  For conic contractions this identifies
their pullback curve fields inside $k(W)$, not merely the abstract
isomorphism classes of the curves.  For identical birational
contractions the common exceptional orbit cancels from the horizontal
cocycle.  Thus the restricted birational map is a structured
isomorphism in either case, with equal vertical divisor data and
$H_Q=0$.

Otherwise the two contractions give distinct $K_W$-negative extremal
rays.  These are the two boundary rays of the two-dimensional cone of
curves of the projective surface $W$.  The anticanonical divisor is
positive on both and is therefore ample.  This proves \textup{(ii)}.
The argument uses numerical classes and contractions over $k$ and
does not require the generic surfaces or conic bases to split.
\end{proof}

\subsection{Finite sets attached to the generic surfaces}
\label{sl:generic-sets}

We now work over the rational surface field $k=\C(Q)$ obtained when
$\dim Q=2$.  Write $\overline k$ for an algebraic closure and
$G_k=\operatorname{Gal}(\overline k/k)$.  All finite $G_k$-sets below carry
continuous actions.  Extend the test to these sets by
\begin{equation}\label{sl:orbit-test}
 u(A)=\sum_{O\in A/G_k}u(k(O)),
\end{equation}
where $k(O)$ is the finite extension corresponding to the transitive orbit
$O$.  These are surface function fields over $\C$.  Blowing up a closed point
with residue field $K$ adds a divisor with field $K(t)$, so its horizontal
contribution is exactly $u(K)$ by Lemma~\ref{pre:mrc-fibration}.
We will repeatedly use
\begin{equation}\label{sl:small-orbits}
 |O|\leq2\quad\Longrightarrow\quad u(k(O))=0.
\end{equation}
This is the small-orbit observation of Lemma~\ref{pre:small-orbits}:
$k$ is rational, and a quadratic extension has a nonidentity
$k$-automorphism, excluded for a test K3 field.

The generic Mori surfaces are rank-one del Pezzo surfaces over $k$, or
relatively rank-one conic bundles $S\to C$ over a geometrically integral
genus-zero curve.  We do not assume that either the surfaces or the curve
$C$ are split over $k$.

\begin{lemma}[Conic fibres and the horizontal vertical correction]
\label{sl:conic-h}
For a conic-bundle generic surface $S\to C$, let $\Delta$ be the finite
$G_k$-set of singular geometric fibres and let $\widetilde\Delta$ be the
double set of their components.  Each component cover over a discriminant
orbit is nonsplit.  The quaternion class of the generic conic over $k(C)$
ramifies exactly at these points, and its residue extensions are exactly
these component covers.  Moreover,
\begin{equation}\label{sl:h-discriminant}
 h_Q(V/B)=u(\widetilde\Delta)-u(\Delta)=-u(\Delta).
\end{equation}
For a rank-one del Pezzo generic surface, $h_Q(V/B)=0$.
\end{lemma}

\begin{proof}
The morphism from the smooth surface to the smooth curve is flat,
so its fibres have the arithmetic genus of the generic conic.
Over $\overline k$, every fibre is connected, has arithmetic genus zero,
and has anticanonical degree two.  The anticanonical degree of each
irreducible component is a positive integer.  An irreducible reduced
fibre therefore has arithmetic genus zero and is a smooth rational curve.
A double fibre $2R$ would have $R^2=0$ and $K_S\cdot R=-1$, contradicting
adjunction parity.  The only remaining possibility is a reduced pair
$R_1+R_2$ with $-K_S\cdot R_i=1$.  Put $m=R_1\cdot R_2$.  Since
$(R_1+R_2)\cdot R_i=0$ and the fibre is connected, adjunction gives
\[
 2p_a(R_i)-2=-m-1,\qquad m\geq1.
\]
Thus $m=1$, both curves are smooth rational curves of self-intersection
$-1$, and they meet transversely.

If the component cover over an orbit were split, the sum of one chosen
component over each point of that orbit would descend to a divisor on
$S$.  It has zero intersection with a general fibre and negative
intersection with every component in that chosen orbit.  This is
impossible when the relative numerical divisor space is one-dimensional:
the general fibre is a nonzero relative curve class and spans its dual.
Consequently every component orbit has field $L$ quadratic over the
corresponding discriminant field $K$.

The relative anticanonical system embeds the fibres as plane conics:
on a smooth fibre it has degree two, and on a reducible fibre it has
degree one on each component, with vanishing first cohomology.  Cohomology
and base change give the local rank-three system.  At a singular fibre
the total surface is regular, so after diagonalizing the quadratic form
over the local discrete valuation ring, its nonsingular binary part has
unit determinant and the remaining coefficient has valuation one.
More explicitly, with uniformizer $t$ the equation is
\[
 ax^2+by^2+tc z^2=0,\qquad a,b,c\text{ units}.
\]
An exponent of $t$ greater than one would make the total space singular
at the geometric node. We use the Kummer identification and quaternion
residue formula of \cite[\S2, equations~(1)--(4)]{AuelBohningBothmerPirutka20}.
The associated quaternion can be written
$(-ab,-act)$; its tame residue is the square class of $-ab$ in the
residue field.  This is exactly the extension separating the two lines
$ax^2+by^2=0$.  At a smooth fibre all three diagonal coefficients are
units, so the residue vanishes.

The corresponding vertical divisor has function field $L(t)$, since a
line is split over its component field.  The involution of $L/K$ excludes
$u(L)=1$ by the automorphism clause of Lemma~\ref{pre:small-orbits};
the degree of $L$ over $k$ need not be two.  All smooth
fibres cancel against their base points by the MRC rule.  This proves
\eqref{sl:h-discriminant}.  On a del Pezzo side the generic base is a
point, so there are no vertical primes dominating $Q$.
\end{proof}

We have computed $h_Q$ for both kinds of generic end.  It remains to
express the horizontal cocycle as an endpoint difference.  The following
finite $G_k$-sets will define a candidate function $w$.  Their orbit fields
matter: equality of rational permutation representations alone need not
identify the fields tested by $u$.

We use the geometric classification, intersection lattices and
anticanonical models of smooth del Pezzo surfaces
\cite[Sections~3.4 and~6.1]{DolgachevIskovskikh2009}.
All plane blowup bases below are chosen over $\overline k$.
For a del Pezzo surface $S$, let
\[
 A(S)=\{F\in\operatorname{Pic}(S_{\overline k}):
           F^2=0,\ -K_S\cdot F=2\}.
\]
For a conic bundle, let $P_\Delta$ be the set of all choices of one
component in each singular geometric fibre.  It has $2^{|\Delta|}$
elements, with its natural $G_k$-action.

The del Pezzo pencil sets and the endpoint-difference calculation below
adapt the virtual N\'eron--Severi viewpoint of Lin, Shinder and Zimmermann
\cite[Definition~5.2 and Proposition~5.5]{LinShinderZimmermann23}.
The conic correction also records the component choices needed for the
generic Mori surfaces here.

The two-ray calculation below will show why only the
degree-five and degree-six pencil sets, and the component choices
over three singular conic fibres, enter this function.  Once that
identity is proved, \eqref{sl:horizontal-residual} will express the
remaining link contribution as the change of $w-h_Q$.

\begin{definition}\label{sl:w-def}
For a generic Mori surface with its specified structure, set
\[
 w(S)=
 \begin{cases}
 u(A(S)),&S\text{ is a rank-one del Pezzo of degree }5\text{ or }6,\\
 \tfrac12u(P_\Delta),&S\to C\text{ is a conic bundle and }|\Delta|=3,\\
 0,&\text{otherwise}.
 \end{cases}
\]
The del Pezzo value depends on the surface; the conic value includes its
chosen genus-zero curve structure.
\end{definition}

For a fourfold node $V\to B$ with $\dim B=2$, write
$w(V/B):=w(V_\eta)$, where $\eta=\Spec\C(B)$.

The factor $1/2$ is harmless even integrally.  Complementing all choices
is a fixed-point-free, equivariant involution of $P_\Delta$.  An orbit
preserved by this involution has a nonidentity field automorphism and
contributes zero.  The remaining orbits occur in isomorphic pairs.
Thus $u(P_\Delta)$ is even.

\begin{lemma}[The pencil sets]\label{sl:pencil-sets}
The sets $A(S)$ have respectively five and three elements in degrees
five and six.  In degree six there is also a two-element set
\[
 A_2(S)=\{C\in\operatorname{Pic}(S_{\overline k}):
                  C^2=1,\ -K_S\cdot C=3\},
\]
and an equality of rational $G_k$-representations
\begin{equation}\label{sl:degree-six-representation}
 \operatorname{Pic}(S_{\overline k})_\Q\oplus\Q
       \simeq\Q[A(S)]\oplus\Q[A_2(S)].
\end{equation}
If $S$ has Picard rank one over $k$, both finite sets in
\eqref{sl:degree-six-representation} are transitive.
\end{lemma}

\begin{proof}
Use a plane blowup basis $L,E_1,\ldots,E_n$ with $n=9-K_S^2$.
A class $xL-\sum y_iE_i$ belongs to $A(S)$ exactly when
\[
 \sum y_i=3x-2,\qquad \sum y_i^2=x^2.
\]
Cauchy's inequality gives $(3x-2)^2\leq nx^2$.  For $n=4$ this restricts
the integer $x$ to $1,2$.  At $x=1$ there is one coefficient $y_i=1$
and all others vanish; at $x=2$ all four coefficients are one.  These
are the four line pencils and the conic pencil.  For $n=3$ the inequality
forces $x=1$, giving the three line pencils.  The same calculation for
$A_2(S)$ gives $x=1,2$ and the two classes
\[
 L,\qquad 2L-E_1-E_2-E_3.
\]
The three pencil classes $L-E_i$ and these two plane classes span the
degree-six Picard space.  Their only linear relation is the equality
of the two class sums, each equal to $-K_S$.  The relation is invariant
under $G_k$; semisimplicity yields \eqref{sl:degree-six-representation}.
Taking invariants gives
\[
 2=\#(A(S)/G_k)+\#(A_2(S)/G_k).
\]
The Brauer obstruction to descending an invariant line bundle is torsion
\cite[Remark~2.7]{AuelBernardara2018}. Thus invariant rational Picard
classes descend, and the invariant Picard space has dimension one.  Both finite
sets are nonempty, proving transitivity.
\end{proof}

\subsection{The two-ray surface calculation}
\label{sl:two-ray-calculation}

\begin{proposition}[Surface-link identity]\label{sl:surface-identity}
Let $\varphi:V/B\dashrightarrow V'/B'$ be an ordinary fourfold link over a
two-dimensional common base $Q$.  Put $k=\C(Q)$, and denote its
generic end surfaces, with their induced Mori structures, by $S$ and
$S'$.  If its restriction is the nontrivial two-ray move in
Lemma~\ref{sl:generic-dichotomy}\textup{(ii)}, then
\begin{equation}\label{sl:surface-equation}
 c_{u,\mathrm{hor}/Q}(\varphi)=w(S')-w(S).
\end{equation}
Every conic surface participating in such a move has at most seven
singular geometric fibres.  In the other case of
Lemma~\ref{sl:generic-dichotomy}, $H_Q=0$ without any bound on the
discriminant size.
\end{proposition}

\begin{proof}
By Lemma~\ref{sl:generic-dichotomy}, the common smooth surface $W$ is
del Pezzo with a two-dimensional invariant numerical plane and the
two end contractions generate its two rays.  Each is either divisorial
to a smooth surface or a conic contraction.  For a divisorial
contraction, pass first to $\overline k$.  Its exceptional intersection
matrix is negative definite.  An exceptional component $R$ has
$R^2<0$ and $K_W\cdot R<0$; adjunction
$2p_a(R)-2=R^2+K_W\cdot R$ forces both intersections to equal $-1$
and $R$ to be a smooth rational curve.  Two such curves must be
disjoint, since otherwise their two-by-two intersection matrix would
not be negative definite.  Thus the contraction is the blowdown of
disjoint $(-1)$-curves to distinct smooth points.  Each Galois orbit
of these curves contributes an independent invariant relative divisor
class.  Relative Picard number one therefore gives a single orbit.
The contraction is consequently the blowup of one closed point over
$k$, and the equivariant correspondence between the geometric points
and their exceptional curves identifies their orbit fields.
Put $D=K_W^2>0$.

Every fibre class used below is the integral class of a geometric
fibre in $\operatorname{Pic}(W_{\overline k})$.  It is $G_k$-invariant,
has square zero and anticanonical degree two, and lies in the rational
numerical plane over $k$ by the torsion descent argument used in
Lemma~\ref{sl:pencil-sets}.  A closed fibre over a point of degree $r$
on the original genus-zero curve has class $r$ times this geometric
class.  Thus none of the following calculations requires a $k$-point
on the conic base.

For a divisorial end let $E$ denote its orbit of exceptional curves,
write $e=|E|$, and put $E_\Sigma=\sum_{R\in E}R$ on
$W_{\overline k}$.  Thus $u(E)=u(k(E))$ counts one orbit field,
whereas $E_\Sigma^2=-e$ counts geometric curves in an intersection
calculation.  When both ends are divisorial, use $F$, $f=|F|$, and
$F_\Sigma$ for the corresponding data at the primed end.  In the diagram
\[
\begin{tikzcd}[column sep=large]
 &W\arrow[dl,"\pi"']\arrow[dr,"\pi'"]&\\
 S\arrow[rr,dashed,"\varphi_k"']&&S'
\end{tikzcd}
\]
$E$ is the exceptional set of $\pi$ and $F$ that of $\pi'$.
The factorization $\varphi_k=\pi'\circ\pi^{-1}$ therefore contributes
$u(E)-u(F)$ by Lemma~\ref{pre:cocycle}: the extraction has positive
sign and the contraction negative sign.  Neither term is multiplied
by its orbit size.

\paragraph{One divisorial end.}
Let $\pi:W\to S$ be the contraction to the rank-one del Pezzo end,
of degree $a$, and let $W\to C$ be the other, conic contraction.
Then $D=a-e>0$.  Write
$H=\pi^*(-K_S)$.  The plane spanned by $H$ and $E_\Sigma$ has
intersection form $H^2=a$, $H\cdot E_\Sigma=0$, and
$E_\Sigma^2=-e$.  A nonzero isotropic fibre class in this rational
plane therefore gives $ax^2=ey^2$ with nonzero rational $x,y$, so
$ae$ is a square.  Since $1\leq e<a\leq9$, the
complete list is
\begin{equation}\label{sl:div-conic-table}
 (a,e)=(4,1),\ (8,2),\ (9,1),\ (9,4).
\end{equation}
Indeed, two positive integers with square product have the same
squarefree part.  Below ten the only unequal such pairs with smaller
entry less than the larger are $1,4$, $2,8$, $1,9$, and $4,9$.

Over $\overline k$, contracting one line in every singular conic fibre
gives a ruled surface over $\mathbf P^1$.  Its canonical square is eight,
so
\begin{equation}\label{sl:conic-size}
 |\Delta|=8-D.
\end{equation}
The first three cases of \eqref{sl:div-conic-table} have extraction
degree at most two and discriminant size different from three.  Their
two $w$ values and their cocycle contributions vanish.

In the remaining case $(a,e)=(9,4)$, work geometrically on the plane
blown up at four points.  Galois fixes $L$ and permutes the four
exceptional classes transitively.  Thus the invariant fibre class
has the form $F=xL-y\sum_{i=1}^4E_i$ with $x,y\in\Z$.  The equations
$F^2=0$ and $-K_W\cdot F=2$ give $x=2$, $y=1$.
The conic pencil is therefore the pencil through the four points.
Its three reducible members correspond to the partitions of these
points into two pairs.  For each point $p$, select in every partition
the pair containing $p$.  This gives a star of three lines; selecting
the other pair gives its complementary triangle.  The four stars
and four triangles exhaust the eight possible selections.  The
construction commutes with every permutation of the four points,
so there is an isomorphism of $G_k$-sets
\begin{equation}\label{sl:eight-orientations}
 P_\Delta\simeq E\sqcup E.
\end{equation}
Consequently $w(W\to C)=u(E)$, which is the extraction contribution.
Reversing the link reverses both signs.

\paragraph{Two conic ends.}
The same surface has two fibre classes $F_1,F_2$.  Put $m=F_1\cdot F_2$.
They are distinct isotropic rays, so $m>0$, and their anticanonical
degrees are two.  In their numerical plane,
\[
 -K_W=\frac2m(F_1+F_2),\qquad
 D=\frac8m,\qquad F_1+F_2=\frac4D(-K_W).
\]
Since $m$ is a positive integer, the equality $D=8/m$ excludes degree five.
Thus neither conic bundle has discriminant size three, by
\eqref{sl:conic-size}.  Their $w$ values vanish and there is no
divisorial contribution.

\paragraph{Two divisorial ends: the numerical list.}
Now let $e,f$ be the orbit sizes, so that the end degrees are $D+e$
and $D+f$, each at most nine.  If $D=1$ or $2$, the Bertini or Geiser involution $\iota$ fixes $K_W$
and acts as $-1$ on $K_W^\perp$
\cite[Sections~6.6 and~6.7]{DolgachevIskovskikh2009}.
It is defined over $k$: the corresponding complete anticanonical or
bi-anticanonical system is defined over $k$, and the nontrivial
involution of its quadratic function-field extension extends to $W$.
If $\iota$ preserved an exceptional ray with orbit sum $E_\Sigma$, then
$\iota_*E_\Sigma=\lambda E_\Sigma$ for some $\lambda>0$; intersection with $-K_W$
forces $\lambda=1$. The fixed-space formula would then give
$E_\Sigma\in\Q K_W$, contrary to $E_\Sigma^2<0<K_W^2$.
Thus $\iota$ exchanges the two contractions and identifies their entire
exceptional $G_k$-sets and output surfaces over $k$. Both sides of
\eqref{sl:surface-equation} vanish.

For the remaining degrees we first restrict the possible orbit sizes by
intersection theory, then identify the finite $G_k$-sets.  The target is
\[
 u(E)-u(F)=w(S')-w(S).
\]
Small orbits contribute zero; the other terms will be matched through
equivariant bijections of exceptional or pencil sets.

Suppose $D\geq3$ and put $m=E_\Sigma\cdot F_\Sigma$.  The Gram determinant of
$-K_W,E_\Sigma,F_\Sigma$, which lie in a plane, gives
\[
 Dm^2-2efm-ef(D+e+f)=0.
\]
Since the two effective exceptional sums meet nonnegatively, the
solution is
\begin{equation}\label{sl:intersection-formula}
 m=\frac{ef+\sqrt{ef(D+e)(D+f)}}D.
\end{equation}
Each curve in an orbit has the same total incidence with the opposite
orbit.  Thus $m/e$ and $m/f$ are integers.  The anticanonical embedding
realizes $(-1)$ curves as lines; two different such curves meet at most
once.  In particular $m\leq ef$.

These conditions give precisely the following necessary possibilities,
with $e\leq f$:
\begin{center}
\begin{tabular}{@{}cl@{}}
\toprule
$D$&$(e,f,m)$\\
\midrule
3&$(2,5,10),\ (3,3,9),\ (6,6,30)$\\
4&$(1,5,5),\ (2,2,4),\ (4,4,12)$\\
5&$(1,3,3)$\\
6&$(3,3,6)$\\
7&$(1,2,2)$\\
\bottomrule
\end{tabular}
\end{center}
Here is a direct exhaustion.  For $e=f$, formula
\eqref{sl:intersection-formula} gives $m/e=1+2e/D$, so $D\mid2e$;
impose also $1\leq e\leq9-D$ and $m\leq e^2$.  This gives exactly
the diagonal entries displayed.  For $e<f$, the values $i(i+D)$ for
$1\leq i\leq9-D$ must have matching squarefree parts.  The lists for
$D=3,4,5,6,7$ are respectively
\[
\begin{array}{c|l}
3&4,10,18,28,40,54\\
4&5,12,21,32,45\\
5&6,14,24,36\\
6&7,16,27\\
7&8,18.
\end{array}
\]
They give only $(2,5),(1,5),(1,3),(1,2)$ in their respective rows.
For $D\geq8$ the permitted range has at most one entry and supplies
neither a diagonal nor an off-diagonal case.

\paragraph{Identifying the finite sets.}
In the cases $(D,e,f)=(3,2,5)$ and $(3,3,3)$, every curve in $F$ meets
every curve in $E$ once.  After contracting $E$, its image class $C$
on $S$ has
\[
 (K_S^2,C^2,-K_S\cdot C)=(5,1,3)
 \quad\text{or}\quad(6,2,4),
\]
respectively.  Therefore $-K_S-C$ is in $A(S)$.  This construction is
injective: equal image classes have equal pullbacks, and all the
multiplicities at the extracted points equal one, so their strict
transforms have the same $(-1)$ class.  A $(-1)$ class has a unique
effective representative.  The cardinalities in Lemma~\ref{sl:pencil-sets}
then show that
\[
 F\simeq A(S)
\]
as actual $G_k$-sets.  For $(3,3,3)$ the construction on the other
side also gives $E\simeq A(S')$.  For $(3,2,5)$ the other side has
degree eight and $u(E)=0$.  Thus the desired difference is $u(E)-u(F)$
in both cases.

For $(D,e,f)=(4,1,5)$ and $(5,1,3)$ every opposing curve meets the
single extracted curve once.  Its image already has square zero and
anticanonical degree two, so the image itself identifies $F$ with
$A(S)$.  The injectivity argument is unchanged.  The opposite end has
degree nine or eight, respectively, and the one-point orbit contributes
zero.  These cases again give \eqref{sl:surface-equation}.

In the diagonal cases $(D,e)=(3,6),(4,4),(6,3)$, the value of $m/e$
is $e-1$.  Each curve therefore misses exactly one curve in the opposite
orbit.  This rule gives an equivariant bijection $E\simeq F$, since
the same statement holds in both directions.  The output degrees are
nine or eight, so no pencil terms occur.  In degree seven both orbit
sizes are at most two and all terms vanish by \eqref{sl:small-orbits}.

It remains to treat $(D,e,f)=(4,2,2)$, whose two outputs have degree
six.  Here we must identify the actual pencil $G_k$-sets, since an
equality of rational permutation representations alone need not
identify their orbit fields.  The Picard blowup decompositions imply
\[
 \operatorname{Pic}(S_{\overline k})_\Q\oplus\Q[E]
 \simeq
 \operatorname{Pic}(S'_{\overline k})_\Q\oplus\Q[F].
\]
Insert \eqref{sl:degree-six-representation} on both sides.  The
permutation representations of $E,F,A_2(S),A_2(S')$ have only
one-dimensional rational constituents, because all four sets have
two elements.  By transitivity, the reduced representation of each
triple $A(S)$ or $A(S')$ is irreducible over $\Q$: its image is either
$C_3$, acting through $\Q(\zeta_3)$, or $S_3$, acting through its
standard representation.  Semisimplicity consequently identifies
these two reduced representations.  Both are faithful on their
image, so their kernels in $G_k$ agree.  The quotient is $C_3$ or
$S_3$; its transitive action on three points is unique up to
isomorphism (the stabilizer is trivial or an order-two subgroup,
and all such subgroups of $S_3$ are conjugate).  Hence
\[
 A(S)\simeq A(S')
\]
as $G_k$-sets.  The two extraction orbits contribute zero by
\eqref{sl:small-orbits}, completing the final case.

Finally, positivity of $D$ in every nontrivial conic case and
\eqref{sl:conic-size} give $|\Delta|\leq7$.  In the structured-isomorphism
case of Lemma~\ref{sl:generic-dichotomy}, the horizontal cocycle is
zero and the $h_Q$ values agree; this case includes cancellation of
identical generic end contractions.  Thus $H_Q=0$ even when the conic
discriminant is arbitrarily large.
\end{proof}

\begin{remark}\label{sl:local-correction}
Combining Proposition~\ref{sl:surface-identity} with
Proposition~\ref{sl:vertical}, the residual on a nontrivial generic
surface link is the change of $w-h_Q$.  On a del Pezzo side this is
$w$; on a conic side it is
\[
 u(\Delta)+\mathbf1_{|\Delta|=3}\,\tfrac12u(P_\Delta).
\]
The surface field $k$ is part of this formula.  Its compatibility
across different two-dimensional common bases is proved in the next
section.  The local calculation alone does not assert that this is
already an invariant of the entire conic Mori fibre space.
\end{remark}

This interpretation agrees with the virtual N\'eron--Severi sets
of Lin--Shinder--Zimmermann
\cite[Proposition~4.8, Definition~5.2 and Proposition~5.5]{LinShinderZimmermann23}:
their point and two-element terms vanish under our test.  The explicit
proof was included to retain the actual orbit fields at every step.
